%%%%%%%%%%%%%%%%%%%%%%%%%%%%%%%%%%%%%%%%%%%%%%%%%%%%%

\documentclass[letterpaper, 10 pt, conference]{ieeeconf}  % Comment this line out if you need a4paper

\IEEEoverridecommandlockouts                        	% This command is only needed if 
                                                         	 		% you want to use the \thanks command
\overrideIEEEmargins                                      	% Needed to meet printer requirements.

% Packages:
\usepackage{graphics} % for pdf, bitmapped graphics files
\usepackage{epsfig} % for postscript graphics files
%\usepackage{mathptmx} % assumes new font selection scheme installed
\usepackage{times} % assumes new font selection scheme installed
\usepackage{amsmath} % assumes amsmath package installed
%\usepackage{amsthm}
\usepackage{amssymb}  % assumes amsmath package installed
\usepackage{listings}
\usepackage{stmaryrd}
\usepackage{graphicx, float, wrapfig}
%\usepackage{subfigure}
\usepackage{subfig}
\usepackage{tikz}
%\usepackage{algpseudocode}
\usepackage{algorithm}
\usepackage{algorithmic}
\usepackage{mathtools}
\usepackage{cite}
\usepackage[font = small]{caption}

% New commands:
\newcommand{\empha}[1]{ \textbf{\emph{#1}}}
\DeclareMathOperator{\F}{\rotatebox[origin=c]{45}{$\Box$}}
\DeclareMathOperator{\X}{\bigcirc}
\DeclareMathOperator{\G}{\Box}
\newcommand{\LTLG}{\G}
\newcommand{\LTLF}{\F}
\newcommand{\LTLX}{\X}
%\newcommand{\NN}{\mathbb{N}}
%\newcommand{\RR}{\mathbb{R}}
%\DeclareMathOperator{\Cox}{\scalebox{0.8}{$\F$}\hspace{-9.3pt}\bigcirc}
\DeclareMathOperator{\Cox}{
    \begin{tikzpicture}[baseline=(X.base)]
    \path[use as bounding box] (-0.16,-0.17) -- (0.16,0.17);
    \node at (0,0) {\scalebox{0.8}{$\F$}}; \node (X) at(0,0) {$\bigcirc$}; 
    %\draw (current bounding box.north west) rectangle (current bounding box.south east);
    \end{tikzpicture}}
    
% Define an "Example" environment with its own counter    
\newcounter{examplecounter}		% Example counter. Independent of theorem and other counters.
\newenvironment{myExample}
{
    	\refstepcounter{examplecounter}%
  	\textbf{Example \arabic{examplecounter}.}% Example #.
  	\quad							% Let's give our text some space :-)
}
    
%%%%%%%%%%%%%%%%%%%%%%%%%%%%%%%%%%%%%%%%%%%%%%%%%%%%%

\title{\LARGE \bf
	Higher Level Abstractions for Provably Correct Robot Planning: % obv, line change does nothing
	Application to Synthesis for Open--Worlds
}

\author{Spyros Maniatopoulos, Matthew Blair and Hadas Kress-Gazit% 	<-this % stops a space
%\thanks{*This work was not supported by any organization}% 			<-this % stops a space
\thanks{The authors are with the Autonomous Systems Lab, School of Mechanical and Aerospace Engineering, Cornell University, Ithaca, NY 14853, USA
        {\tt\small \{sm2296, hadaskg\}@cornell.edu}}%
%\thanks{$^{2}$Bernard D. Researcheris with the Department of Electrical Engineering, Wright State University,
%        Dayton, OH 45435, USA
 %     {\tt\small b.d.researcher@ieee.org}}%
}

%%%%%%%%%%%%%%%%%%%%%%%%%%%%%%%%%%%%%%%%%%%%%%%%%%%%%

\begin{document}

\maketitle
\thispagestyle{empty}
\pagestyle{empty}

%%%%%%%%%%%%%%%%%%%%%%%%%%%%%%%%%%%%%%%%%%%%%%%%%%%%%

\begin{abstract}

Check LTL in LaTeX: $\X \G \F \Cox$.

\end{abstract}

%%%%%%%%%%%%%%%%%%%%%%%%%%%%%%%%%%%%%%%%%%%%%%%%%%%%%

\section{INTRODUCTION}

...

\begin{itemize}
	\item High-level robot mission planning
  	\item Fixed workspace (state-space) assumption
 	\item Recent literature
   	\item Brief statement of our contributions
   	\item Paper structure
\end{itemize}

Some references: \cite{MurrayICRA2012}, \cite{MurrayICRA2013a},  \cite{BeltaICRA2012}, \cite{Dimos2013ICRA}, \cite{Belta2013RSS}, \cite{BingxinRSS2012}.

%%%%%%%%%%%%%%%%%%%%%%%%%%%%%%%%%%%%%%%%%%%%%%%%%%%%%

\section{PRELIMINARIES}

\subsection{Controller Synthesis}

...

\subsection{LTL Mission Planning}

LTLMoP, Structured English, Parsing

%%%%%%%%%%%%%%%%%%%%%%%%%%%%%%%%%%%%%%%%%%%%%%%%%%%%%

\section{HIGHER LEVEL ABSTRACTIONS}

\subsection{Groups of Propositions} 

...

\begin{myExample}\label{Ex:groups}
	Groups of sensors or actions (use Hadas's and Cam's ideas).
\end{myExample}

\subsection{Group Quantifiers} 

...

\begin{myExample}\label{Ex:quantifiers}
	2-3 one-liners that demonstrate the 3 quantifiers (any, all, each)
\end{myExample}

\subsection{Proposition Correspondence} 

...

\begin{myExample}\label{Ex:corresponding}
	Robowaiter (orders-customers) or PR3 (letters-addresses)
\end{myExample}

%%%%%%%%%%%%%%%%%%%%%%%%%%%%%%%%%%%%%%%%%%%%%%%%%%%%%

\section{SYNTHESIS FOR OPEN--WORLDS}

\subsection{Tasks over Unknown Propositions}

...

% \subsection{Exploring the Physical Workspace}

\subsection{Augmenting the State-Space} % Re-synthesis (not local yet)

...

%%%%%%%%%%%%%%%%%%%%%%%%%%%%%%%%%%%%%%%%%%%%%%%%%%%%%

\section{SIMULATION: Open--World Exploration}  % or make exploration a subsection in pref section

...

%%%%%%%%%%%%%%%%%%%%%%%%%%%%%%%%%%%%%%%%%%%%%%%%%%%%%

\section{CONCLUSIONS AND FUTURE WORK}

...

Local/Incremental re-synthesis

%\addtolength{\textheight}{-12cm}   % This command serves to balance the column lengths
                                  % on the last page of the document manually. It shortens
                                  % the textheight of the last page by a suitable amount.
                                  % This command does not take effect until the next page
                                  % so it should come on the page before the last. Make
                                  % sure that you do not shorten the textheight too much.

%%%%%%%%%%%%%%%%%%%%%%%%%%%%%%%%%%%%%%%%%%%%%%%%%%%%%

\section*{APPENDIX}

...

\section*{ACKNOWLEDGMENT}

... Cameron ...

%%%%%%%%%%%%%%%%%%%%%%%%%%%%%%%%%%%%%%%%%%%%%%%%%%%%%

\bibliographystyle{ieeetran}
\bibliography{../../common/spyros-biblio.bib}

\end{document}
